% ============================================================
% MODULO 02 - RESUMO ROTULADO (padrao do modelo da disciplina)
% NBR 6028:2021 — resumo em paragrafo unico, 3a pessoa, voz ativa
% ============================================================
\begin{center}
{\bfseries RESUMO}
\end{center}

\noindent\textbf{Introdu\c{c}\~ao:} o Transtorno de D\'eficit de Aten\c{c}\~ao e Hiperatividade (TDAH) de apresenta\c{c}\~ao combinada em pr\'e-escolares compromete a perman\^encia em atividades, o v\'inculo terap\^eutico e a conviv\^encia entre pares, e nessa faixa et\'aria a evid\^encia farmacol\'ogica \'e limitada. O caso que motivou este trabalho chegou \`a cl\'inica humanista assim: M., 5 anos, laudo de TDAH misto, n\'ivel de hiperatividade grave, a ponto de n\~ao conseguir permanecer diante da televis\~ao. Depois que a fam\'ilia mudou de endere\c{c}o, a crian\c{c}a come\c{c}ou a brincar todos os dias com outras crian\c{c}as da vizinhan\c{c}a; a m\~ae e a terapeuta passaram a not\'a-la bem mais tranquila, inclusive durante as sess\~oes. \textbf{Objetivos:} reunir e discutir a literatura publicada entre 2004 e 2026 sobre o brincar livre e social em espa\c{c}os externos do cotidiano -- rua, quintal, vizinhan\c{c}a, \'areas verdes -- e sua rela\c{c}\~ao com a intensidade dos sintomas do TDAH, com o prop\'osito de ler a mudan\c{c}a observada \`a luz do que j\'a se conhece; em particular, mapear as evid\^encias sobre natureza e gravidade do transtorno, sobre o brincar entre pares e a autorregula\c{c}\~ao de crian\c{c}as pequenas, e sobre o lugar da cl\'inica humanista nesse processo. \textbf{Metodologia:} revis\~ao integrativa de car\'ater narrativo, conduzida nas bases SciELO, BVS/PEPSIC, Portal de Peri\'odicos CAPES, PubMed e PsycINFO, com descritores DeCS/MeSH (TDAH; pr\'e-escolar; brincar; meio ambiente; terapia centrada no cliente). O caso foi tratado como vinheta anonimizada, com Termo de Consentimento Livre e Esclarecido do respons\'avel e assentimento l\'udico [A ANEXAR], em conformidade com o Estatuto da Crian\c{c}a e do Adolescente (Lei n. 8.069/1990), a Lei Geral de Prote\c{c}\~ao de Dados (Lei n. 13.709/2018) e as Resolu\c{c}\~oes n. 466/2012 e n. 510/2016 do Conselho Nacional de Sa\'ude. A fam\'ilia n\~ao informou uso de medicamento at\'e o fechamento deste trabalho, o que se registra literalmente. \textbf{Resultados e discuss\~oes:} seis refer\^encias-\^ancora, organizadas em tr\^es eixos -- verde e restaura\c{c}\~ao atencional; brincar entre pares como regulador; cl\'inica humanista como ambiente facilitador --, convergem para um mesmo achado: o tipo de TDAH permanece, mas a intensidade dos sintomas diminui, com mais calma, capacidade de esperar a vez e abertura ao contato social. O percurso de M. espelha o relatado por Damasceno et al. (2025) e se mostra compat\'ivel com Kuo e Faber Taylor (2004) e Hood e Baumann (2024); ainda assim, sem escala padronizada de linha de base, sem medida da dose de brincar e sem informa\c{c}\~ao farmacol\'ogica, n\~ao cabe concluir efic\'acia, apenas coer\^encia te\'orica. \textbf{Conclus\~ao:} a vinheta dialoga com a literatura, sem valor probat\'orio; a ecologia do brincar merece lugar na formula\c{c}\~ao cl\'inica do TDAH pr\'e-escolar, e a agenda de continuidade inclui SNAP-IV, registro do tempo de brincar ao ar livre e controle de vari\'aveis de confus\~ao.

\vspace{0.75cm}

\noindent\textbf{Palavras-chave:} Transtorno do D\'eficit de Aten\c{c}\~ao com Hiperatividade; Pr\'e-escolar; Brincar Livre; Psicologia Humanista; Relato de Experi\^encia.

\vspace{1.0cm}

\noindent\textbf{Nota de integridade:} este manuscrito contou com apoio de intelig\^encia artificial na fase de estrutura\c{c}\~ao e foi revisto integralmente pelas autoras, que assumem a responsabilidade pelo conte\'udo final. Nenhum participante, n\'umero, DOI ou aprova\c{c}\~ao \'etica foi inventado; campos pendentes aparecem como [A PREENCHER]/[A ANEXAR]. As refer\^encias centrais foram conferidas em fonte prim\'aria em 5 out. 2026. Recomenda-se a verifica\c{c}\~ao de similaridade em ferramenta institucional antes da submiss\~ao. Conflito de interesses: nada a declarar. Financiamento: nenhum. A classifica\c{c}\~ao Qualis e o aceite dependem do peri\'odico e da banca.

\thispagestyle{empty}
\newpage
\pagenumbering{arabic}
\setcounter{page}{1}
